\documentclass[11pt]{article}
\usepackage{amsmath,amsfonts}
\usepackage{graphicx}
%\usepackage{xcolor}
%\definecolor{gray}{rgb}{.75,.75,.75}
%\usepackage[landscape]{geometry}
%\usepackage{hyperref}
%\usepackage[bookmarks=true, pdfborder={0 0 .5}, urlbordercolor=gray]{hyperref}

\renewcommand{\thefootnote}{\fnsymbol{footnote}}


\addtolength{\topmargin}{-0.4in}
\addtolength{\textheight}{0.8in}
\addtolength{\oddsidemargin}{-.4in}
\addtolength{\textwidth}{.8in}
\pagestyle{empty}

\newcommand{\ds}{\displaystyle}
\newcommand{\ts}{\textstyle}

\newcommand{\Z}{\mathbb{Z}}
\newcommand{\R}{\mathbb{R}}
\newcommand{\C}{\mathbb{C}}
\newcommand{\EE}{\mathcal{E}}
\newcommand{\tZ}{\tilde{\Z}}

\newcommand{\Sym}{\mathrm{Sym}}

\renewcommand{\circle}[1]{{\bigcirc\hspace{-.75em}#1}}

\begin{document}
\centerline{\bf \Large\bf
Geometric Trig Substitution}
\vspace{1em}

\noindent
Here is some esoteric knowledge for the hard-core students.
Consider any integral in which trig functions are combined 
by the four arithmetic operations, such as:
$$
\int \!\!\frac{\cos^2(x)\sin(x)-2\tan(x)+1}{\sec^3(\theta)+\sin^3(\theta)+3\cos(\theta)\sin(\theta)+5}\ d\theta.
$$
There is an amazing technique, the {\it Tangent Half-Angle Substitution}, 
which allows us to reduce any such problem to
the integral of a rational function (a quotient of polynomials), 
which can then be done  by Partial Fractions (see \S7.4).

This substitution is motivated by the geometry of the circle.
Recall that the basic trig functions are {\it circular functions},
meaning we can trace the points on a unit circle by
the coordinates $(x,y)=(\cos(\theta),\sin(\theta))$.
There is another way to trace this circle, the {\it rational parametrization}: for any number $t\in(-\infty,\infty)$,
draw the line $L$ from the fixed point $(-1,0)$ to the point $(0,t)$ on the $y$-axis:
this line cuts the circle in exactly one other point $(x,y)$.
As $(0,t)$ moves along the $y$-axis, the point $(x,y)$ moves around the
entire circle, leaving out only the fixed point $(-1,0)$.
\\[.7em]
\centerline{
\includegraphics[width=3.5in]
{7-3-Rat-Param-2.png} 
}
\vspace{-.9em}

\noindent
To compute the coordinates $(x,y)$ corresponding to a given $t$, 
we write the slope of line $L$ using the two similar
triangles between $L$ and the $x$-axis:
$$
\text{slope} \ =\ \frac{t}{1}\ =\ \frac{y}{x{+}1} 
\qquad\Longrightarrow\qquad
y=t(x{+}1).
$$
The point $(x,y)$ also satisfies the circle equation $x^2+y^2=1$.
Substituting for $y$ gives:
$$
x^2 + (t(x{+}1))^2 = 1
\qquad\Longrightarrow\qquad
x^2 + t^2x^2 + 2t^2x + t^2 - 1= 0\,.
$$
Now, for any $t$, this equation is always satisfied by the fixed point 
with $x=-1$, so $x{+}1$ must be a factor of the last polynomial.
Long division gives:
$$
x^2 + t^2x^2 + 2tx + t - 1
\ =\
(x{+}1) (t^2 x{+}x{+}t^2{-}1) = 0
\quad\Longleftrightarrow\quad
\left\{\begin{array}{@{\!}c}
x = \dfrac{1-t^2}{1+t^2}
\\[1em]
 \text{ or } \ x = -1 \,.
\end{array}\right.$$
To write $y$ in terms of $t$, plug this into $y = t(x+1)$, 
Hence we get a new formula tracing the points of the circle,
controlled by $t\in(-\infty,\infty)$:
$$
(x,y)  = \left(\frac{1-t^2}{1+t^2}\,, \frac{2t}{1+t^2}\right)\!.
$$
Incidentally, plugging in rational values of $t$ produces infinitely
many points on the unit circle with rational coordinates, corresponding
to all Pythagorean triples, whole numbers $(a,b,c)$ with $a^2+b^2=c^2$.
For example, $t = 2$ gives:
$$
(x,y)  = \left(-\tfrac{3}{5}, \tfrac{4}{5}\right)
\quad\Longrightarrow\quad
(\tfrac{3}{5})^2+(\tfrac{4}{5})^2=1
\quad\Longrightarrow\quad
3^2+4^2=5^2.
$$

Now, comparing $(x,y)  = \left(\frac{1-t^2}{1+t^2}\,, \frac{2t}{1+t^2}\right)$ 
with  $(x,y)=(\cos(\theta),\sin(\theta))$ suggests the substitution:
$$
\cos(\theta) = \frac{1-t^2}{1+t^2}, \qquad
\sin(\theta) =  \frac{2t}{1+t^2}, \qquad
\tan(\theta) = \frac{2t}{1-t^2}, \qquad \text{etc.}
$$
We can think of this as a backward substitution (\S7.3),
$\theta = \arcsin\!\left(\!\frac{2t}{1+t^2}\!\right)$,
and compute:
$$
\sin(\theta) =  \frac{2t}{1+t^2}
\quad\Longrightarrow\quad
\cos(\theta)\,d\theta \ =\ 2\,\frac{1-t^2}{(1+t^2)^2}\,dt
$$
$$
\quad\Longrightarrow\quad
d\theta \ =\ 2\,\frac{1}{\cos(\theta)}\,\frac{1-t^2}{(1+t^2)^2}\,dt
\ =\
2\,\frac{1+t^2}{1-t^2}\,\frac{1-t^2}{(1+t^2)^2}\,dt
\ =\
\frac{2}{1+t^2}\,dt\,.
$$
\vspace{.1em}

To restore the original variable $\theta$ at the end of the integration, 
we need to write $t=\frac{y}{x+1}$ in terms of trig functions.  
We can do this just by $(x,y)=(\cos(\theta),\sin(\theta))$;
but also recall the theorem of elementary geometry which says the angle between $L$ and the $x$-axis is $\frac\theta2$, giving another expression for the slope:
$$
\text{slope} \ =\ t \ =\ \frac{y}{x{+}1}
\ =\ 
\frac{\sin(\theta)}{\cos(\theta)+1}
\ =\
\tan\!\left(\frac\theta2\right)
\,,
$$
which is why we call it the Tangent Half-Angle Substitution.

\newpage

\centerline{\sc summary} 

\vspace{.5em}

\boxed{$$ \ds
\ \ \ \cos(\theta) = \frac{1-t^2}{1+t^2}, \ \  \ 
 \sin(\theta) =\frac{2t}{1+t^2}, \ \ \ 
d\theta = \frac{2}{1+t^2}\,dt, \ \ \ 
t = \frac{\sin(\theta)}{\cos(\theta)+1}
= \tan\!\left(\frac\theta2\right). \ \ \ 
$$}

\vspace{1.5em}

\noindent
{\sc example:} We carry out this substitution, then the
Partial Fraction Method from \S7.4:
$$
\int\!\!\frac{1}{\sin^2(\theta) + \cos(\theta) + 2}\ d\theta
\ =\ 
\int\!\!\frac{1} { \left(\!\frac{2t}{1{+}t^2}\!\right)^{\!\!2} + \left(\!\frac{1{-}t^2}{1{+}t^2}\!\right) + 2}\ 
\cdot \frac{2}{1{+}t^2}\,dt
\ =\
\int\!\! \frac{2(t^2{+}1)}{t^4 + 8t^2 + 3}\,dt
$$
$$
\ =\
\int\!\! \frac{2(t^2{+}1)}{(t^2{+}4)^2-13}\,dt
\ =\
\int  
\frac {1+\frac3{\sqrt{13}}} {t^2{+}4{+}\sqrt{13}}
\ +\
\frac {1-\frac3{\sqrt{13}}} {t^2{+}4{-}\sqrt{13}}
\ \,dt
$$
\\[-1em]
$$
\ =\ 
(\sqrt{13}{+}3){\ts\sqrt{4{-}\sqrt{13}}}\arctan\!\left(\!\tfrac{t}{\sqrt{4+\sqrt{13}}}\right)
\ +\
(\sqrt{13}{-}3){\ts\sqrt{4{+}\sqrt{13}}}\arctan\!\left(\!\tfrac{t}{\sqrt{4-\sqrt{13}}}\right)
$$
$$
\ =\ 
(\sqrt{13}{+}3){\ts\sqrt{4{-}\sqrt{13}}}\arctan\!\left(\!\tfrac{\tan(\theta/2)}{\sqrt{4+\sqrt{13}}}\right)
\ +\
(\sqrt{13}{-}3){\ts\sqrt{4{+}\sqrt{13}}}\arctan\!\left(\!\tfrac {\tan(\theta/2)}{\sqrt{4-\sqrt{13}}}\right).\!\!
$$
\\[-1.5em]

The point here is not the specific answer, which can be gotten by computer
much more reliably than by hand. 
It is the principle that this is possible for any integral of this type, 
precisely because of the rational parametrization of the circle. 
This leads toward the theory of Lie groups, which generalizes the circle 
to highly symmetric geometric objects in higher dimensions,
starting with the 3-dimensional sphere.\footnote{
We should be able to produce a similar substitution for
integrals involving matrix coefficients of any representation of a compact
Lie group, composed with the exponential map on the Lie algebra.
}


\end{document}

%%%%%%%%%%%  Picture  %%%%%%%%%%
\\[0em]
\centerline{
%\includegraphics[width=2in]
{1-8-Discont-iv.png}
\qquad \qquad 
%\includegraphics[width=2in]
{1-8-Discont-v.png} 
}
\vspace{0em}

\noindent
%%%%%%%%%%%%%%%%%%%%%%%%%

%%%%%%%%%%%  Table  %%%%%%%%%%
In fact, we get the following derivatives:
$$\begin{array}{|c||c|c|c|c|c|c|}
\hline &&&&&& \\[-.9em] 
f(x) & \sin(x) & \cos(x) & \tan(x)  & \sec(x) & \csc(x) & \cot(x)
\\[.2em] \hline &&&&&& \\[-.9em]
f'(x) &\cos(x) & -\sin(x) & \sec^2(x) & \tan(x)\sec(x) & -\cot(x)\csc(x) & -\csc^2(x)
\\[.1em] \hline
\end{array}$$
\\[1em]
%%%%%%%%%%%%%%%%%%%%%%%%%%



















